\section*{अध्याय \ref{ch:b1:erosion-transport} --- अपरदन और परिवहन}

\begin{solution}{exo:b1:erosion-transport:1}
\qty{0.4}{g.L^{-1}} यानी \qty{0.4}{kg.m^{-3}}, इसलिए फ़्लक्स $500 \times 0.4 =
\qty{200}{kg.s^{-1}}$ है। एक वर्ष \num{365.25}~दिनों का \qty{3.156e7}{s} है, इसलिए वार्षिक फ़्लक्स
\qty{6.31e9}{kg}, या
\qty{6.31}{Mt}।
\end{solution}

\begin{solution}{exo:b1:erosion-transport:2}
$w_s = R g d^2/(18\nu) = 1.6548 \times 9.81 \times (\num{30e-6})^2/(18 \times
\num{1.0034e-6}) = \qty{8.09e-4}{m.s^{-1}}$, प्रति सेकंड लगभग \qty{0.8}{mm}। उसकी
कण रेनॉल्ड्स संख्या $w_s d/\nu = \num{8.09e-4} \times
\num{30e-6}/\num{1.0034e-6} = 0.024$ है, 1 से बहुत नीचे: स्टोक्स का नियम लागू है। वह
$2/\num{8.09e-4} = \qty{2472}{s}$ में \qty{2}{m} गिरता है, लगभग
\qty{41}{min}।
\end{solution}

\begin{solution}{exo:b1:erosion-transport:3}
$\tau_b = \rho g h S = 998.2 \times 9.81 \times 0.5 \times \num{4e-4} =
\qty{1.96}{Pa}$। $(\rho_s - \rho) g = 1651.8 \times 9.81 =
\qty{16204}{N.m^{-3}}$: \qty{1}{mm} के लिए, $\theta = 1.96/(16204 \times 0.001) =
0.121$, देहली \num{0.0314} से लगभग चार गुना, इसलिए बालू चलती है; \qty{10}{mm} के लिए
$\theta = 0.0121$, देहली \num{0.0556} का पाँचवाँ भाग, इसलिए बजरी रुकती है। धारा अपनी बालू एक ऐसी बजरी-तल पर ढोती है जो नहीं चलता।
\end{solution}

\begin{solution}{exo:b1:erosion-transport:4}
(क)~$12.6/3.9 = 3.2$: नदियाँ घुले की तुलना में निलंबन में लगभग तीन गुना अधिक ढोती हैं।
(ख)~बड़े जलाशय: $0.20 \times 14 = \qty{2.8}{Gt}$ प्रति वर्ष; सारे: $0.26 \times 14 = \qty{3.64}{Gt}$ प्रति वर्ष। (ग)~घुला फ़्लक्स \qty{3.9e15}{g} प्रति वर्ष, \qty{120}{g.m^{-3}} पर, इसलिए उसे ढोने वाला जल
$\num{3.9e15}/120 = \qty{3.25e13}{m^3}$ है, लगभग \qty{32500}{km^3} प्रति वर्ष। उस जल में \omterm{def:b1:erosion-transport:load}{निलंबित भार}
$\qty{1.26e16}{g}/\qty{3.25e16}{L} =
\qty{0.39}{g.L^{-1}}$ है।
\end{solution}

\begin{solution}{exo:b1:erosion-transport:5}
(क)~खंड का क्षेत्रफल $10 \times 2 = \qty{20}{m^2}$ और गीली परिधि $10 + 2 \times 2 = \qty{14}{m}$ है, इसलिए $R_h = 20/14 = \qty{1.43}{m}$।
(ख)~$\tau_b = 998.2 \times 9.81 \times 1.43 \times 0.001 = \qty{14.0}{Pa}$, \qty{2}{m} गहरी बहुत चौड़ी नहर के $\qty{19.6}{Pa}$ के मुकाबले: अनुपात $R_h/h = 0.714$ है, किनारे घर्षण का शेष भाग लेते हैं। (ग)~चौड़ाई $b$ की आयताकार नहर के लिए, $R_h = bh/(b + 2h) = h/(1 + 2h/b)$, जो $h$ की ओर जाता है, जब $b$ बढ़ता है।
\end{solution}

\begin{solution}{exo:b1:erosion-transport:6}
(क)~$d_1 = \left(18\nu^2/(Rg)\right)^{1/3} = (18/(Rg))^{1/3}\,\nu^{2/3}$, और $R$ तथा $g$ श्यानता पर निर्भर नहीं करते। (ख)~$\nu = \num{1.3059e-3}/999.7 =
\qty{1.306e-6}{m^2.s^{-1}}$, \qty{10}{\celsius} पर, \qty{20}{\celsius} के \qty{1.0034e-6}{m^2.s^{-1}} के मुकाबले: अनुपात \num{1.302} है,
और $1.302^{2/3} = 1.192$, इसलिए $d_1 = 0.104 \times 1.192 = \qty{0.124}{mm}$। (ग)~स्टोक्स के क्षेत्र में $w_s$, $\nu$ का व्युत्क्रमानुपाती है: कण $1.302$ गुना धीमे डूबता है, लगभग \qty{30}{\percent}। ठंडा जल बारीक तलछट को निलंबन में अधिक देर रखता है: शीतकाल में एक नदी उसी प्रवाह के लिए अपनी गाद कुछ और दूर ले जाती है।
\end{solution}

\begin{solution}{exo:b1:erosion-transport:7}
$R g/(18\nu) = \qty{8.99e5}{m^{-1}.s^{-1}}$ के साथ, और संयुक्त नियम से:
\begin{center}
\begin{tabular}{@{}lccc@{}}
\hline
$d$ & स्टोक्स (\si{m.s^{-1}}) & संयुक्त (\si{m.s^{-1}}) & अनुपात\\
\hline
\qty{0.1}{mm} & \num{8.99e-3} & \num{7.53e-3} & \num{1.19}\\
\qty{0.5}{mm} & \num{0.225} & \num{0.0711} & \num{3.16}\\
\qty{2}{mm} & \num{3.60} & \num{0.197} & \num{18.3}\\
\hline
\end{tabular}
\end{center}
अनुपात आकार के साथ बढ़ता है, क्योंकि कण रेनॉल्ड्स संख्या बढ़ती है: स्टोक्स का नियम, जो केवल छोटी रेनॉल्ड्स संख्याओं के लिए वैध है, $d^2$ के साथ बढ़ता है, जबकि पीछे की धारा का कर्षण मोटे कणों को उस वेग तक धीमा करता है जो केवल $\sqrt d$ के साथ बढ़ता है।
\end{solution}

\begin{solution}{exo:b1:erosion-transport:8}
(क)~$(Rg/\nu^2)^{1/3} = \qty{25263}{m^{-1}}$। \qty{0.2}{mm} पर, $D_* = 5.05$
और $\theta_c = 0.3/7.06 + 0.055(1 - e^{-0.101}) = 0.0425 + 0.0053 =
0.0478$; $\tau_c = 0.0478 \times 16204 \times \num{2e-4} = \qty{0.155}{Pa}$।
\qty{20}{mm} पर, $D_* = 505$ और $\theta_c = 0.3/607 + 0.055(1 - e^{-10.1}) =
0.0005 + 0.0550 = 0.0555$; $\tau_c = 0.0555 \times 16204 \times 0.02 =
\qty{18.0}{Pa}$। (ख)~$h = \tau_c/(\rho g S)$ के साथ $\rho g S =
\qty{19.58}{Pa.m^{-1}}$: \qty{7.9}{mm} जल बालू को हिलाता है, \qty{0.92}{m}
बजरी को।
\end{solution}

\begin{solution}{exo:b1:erosion-transport:9}
द्रोणी \qty{0.8}{Mt} = \qty{8e5}{t} प्रति वर्ष खोती है, या $\num{8e5}/2000 =
\qty{400}{t.km^{-2}.yr^{-1}}$। \cref{prop:b1:erosion-transport:denudation} से,
$D = \qty{8e8}{kg}/(2650 \times \qty{2e9}{m^2}) = \qty{1.51e-4}{m}$ प्रति वर्ष, या
\qty{151}{m/Myr}। यह बेरिलियम-10 से मापी गई द्रोणियों की माध्यिका \qty{54}{m/Myr} का लगभग तिगुना है, फिर भी उनके माध्य \qty{218}{m/Myr} से नीचे: एक काफ़ी तेज़ द्रोणी, जैसे तीखे या विवर्तनिक रूप से सक्रिय देश में।
\end{solution}

\begin{solution}{exo:b1:erosion-transport:10}
(क)~$u_* = \sqrt{g h S} = \sqrt{9.81 \times 3 \times \num{5e-4}} =
\qty{0.1213}{m.s^{-1}}$ और $\tau_b = \rho u_*^2 = \qty{14.7}{Pa}$।
(ख)~निलंबन के लिए $w_s < u_* = \qty{0.1213}{m.s^{-1}}$ चाहिए। संयुक्त नियम \qty{1}{mm} पर \qty{0.126}{m.s^{-1}} देता है; $w_s(d) = 0.1213$ को प्रयोग से हल करने पर $d = \qty{0.94}{mm}$ मिलता है। (ग)~बजरी के लिए $\theta_c \approx 0.0555$, इसलिए $\tau_c = \qty{14.7}{Pa}$ के लिए $d \approx
14.7/(0.0555 \times 16204) = \qty{16.3}{mm}$; उस आकार पर देहली-वक्र सचमुच $\theta_c = 0.0555$ देता है। (घ)~लगभग \qty{16}{mm} से मोटे कण तल पर रहते हैं और एक स्थिर बजरी-फ़र्श बनाते हैं; बारीक बजरी और लगभग एक मिलीमीटर से मोटी बालू \omterm{def:b1:erosion-transport:load}{तल-भार} के रूप में चलती हैं; बारीक बालू, गाद और मृत्तिका निलंबन में चलती हैं।
\end{solution}

\begin{solution}{exo:b1:erosion-transport:11}
(क)~क्षेत्र पर औसतन $0.30 \times \qty{1000}{m} = \qty{300}{m}$ चट्टान। (ख)~प्रत्यावर्तन $300 \times 2700/3300 = \qty{245}{m}$ है। (ग)~माध्य सतह केवल $300 - 245 = \qty{55}{m}$ नीची हुई है। (घ)~शिखर, जो अपरदित नहीं हुए, पूरे क्षेत्र के साथ ऊपर उठाए जाते हैं: अब वे पहले से \qty{245}{m} ऊँचे खड़े हैं। घाटियाँ काटने से एक श्रेणी के शिखर उठ सकते हैं।
\end{solution}

\begin{solution}{exo:b1:erosion-transport:12}
(क)~$0.9 \times \qty{5e8}{kg}/\qty{1300}{kg.m^{-3}} = \qty{346000}{m^3}$ प्रति वर्ष,
या जलाशय का \qty{0.69}{\percent}। (ख)~जलाशय का आधा, \qty{25e6}{m^3}, इस दर पर $\num{25e6}/\num{346000} = \num{72}\text{ वर्ष}$ में भरता है। (ग)~जैसे-जैसे जलाशय भरता है, उसका आयतन, और इसलिए जल का आवास-काल, घटता है, और बारीक कण गुज़रते हैं: उसकी फँसाने की दक्षता गिरती है। इस बीच पुराने निक्षेप दबते और संहत होते हैं, पानी निकालते हैं, ताकि उनका थोक घनत्व बढ़े और हर टन कम जगह ले। दोनों प्रभाव आयतन की हानि को जलाशय के पुराने होने पर धीमा करते हैं।
\end{solution}

\begin{solution}{pb:b1:erosion-transport:1}
\textbf{1.} एक वर्ष के जल का आयतन $1.2 \times \num{3.156e7} = \qty{3.79e7}{m^3}$ है;
\qty{8e7}{m^2}
पर फैला, एक \qty{0.473}{m} परत: \qty{473}{mm} प्रति वर्ष का अपवाह।

\textbf{2.} $1.2 \times \qty{0.150}{kg.m^{-3}} = \qty{0.18}{kg.s^{-1}}$, या
\qty{5680}{t} प्रति वर्ष।

\textbf{3.} $1.2 \times 0.080 = \qty{0.096}{kg.s^{-1}}$, या \qty{3030}{t} प्रति
वर्ष।

\textbf{4.} \omterm{def:b1:erosion-transport:load}{तल-भार} \qty{568}{t} प्रति वर्ष है, और कुल $5680 + 3030 +
568 = \qty{9278}{t}$ प्रति वर्ष।

\textbf{5.} निलंबित तलछट की सांद्रता बाढ़ के कुछ दिनों में ऊँची होती है, जब निस्सरण भी ऊँचा होता है, और बाकी समय कम। एक समय-माध्य हर दिन को समान भार देता है, इसलिए वह कम निस्सरण और कम सांद्रता के बहुत दिनों से प्रभावित होता है। \cref{prop:b1:erosion-transport:flux} से वार्षिक द्रव्यमान $\int Q\,C\,\mathrm dt$ है; प्रवाह-भारित माध्य यह समाकल $\int Q\,\mathrm dt$ से भाग देने पर है, इसलिए वही प्रवाह-भारित माध्य है जिसे माध्य निस्सरण से गुणा करना चाहिए।

\textbf{6.} $3030/9278 = \num{0.33}$, एक तिहाई। विश्व की नदियों के लिए,
$3.9/(12.6 + 3.9) = \num{0.24}$, और यदि \omterm{def:b1:erosion-transport:load}{तल-भार} गिना जाता तो और कम। द्रोणी का भाग अधिक है: नीची पहाड़ियाँ, जहाँ धीमा यांत्रिक अपरदन चट्टान को घुलने का समय देता है, और चूना-पत्थर तथा खड़िया, जो पूरी तरह घुलते हैं।

\textbf{7.} $9278/80 = \qty{116}{t.km^{-2}.yr^{-1}}$।

\textbf{8.} $D = \qty{9.278e6}{kg}/(2700 \times \qty{8e7}{m^2}) =
\qty{4.30e-5}{m}$ प्रति वर्ष, \qty{43.0}{m/Myr}।

\textbf{9.} $43.0/54 = 0.80$: विमापित द्रोणियों की माध्यिका से थोड़ा नीचे, और उनके माध्य से पाँच गुना नीचे। एक मध्यम धीमी द्रोणी, जैसी शांत क्षेत्र की नीची पहाड़ियों से अपेक्षित।

\textbf{10.} गॉजिंग कुछ वर्षों के मापन पर माध्य करती है, और उन बाढ़ों पर निर्भर करती है जो उनमें हुईं। बेरिलियम-10 उन कुछ मीटरों के समय पर माध्य करता है जिनमें वह उत्पन्न होता है: \qty{43}{m/Myr} पर, हर मीटर \num{23000} वर्ष लेता है, इसलिए माध्य दसियों हज़ार वर्षों को समेटता है। वे भिन्न हो सकते हैं, दुर्लभ बड़ी बाढ़ों के कारण जो एक छोटे अभिलेख से छूट जाती हैं, इसलिए कि खेती और सफ़ाई ने पिछली सदियों में अपरदन बढ़ाया है, और इसलिए कि तलछट बाढ़-मैदान में संचित होती है और वहाँ से वापस ली जाती है।

\textbf{11.} घुले भार का एक भाग चट्टान से नहीं आता: बाइकार्बोनेट का बहुत-सा भाग हवा या मिट्टी से लिया कार्बन ढोता है, और वर्षा समुद्री लवण लाती है। उन्हें गिनने से द्रोणी को उस चट्टान का श्रेय मिलता जो उसने नहीं खोई।
\textbf{12.} $\num{1.5e6}/1.2 = \qty{1.25e6}{s}$, \num{14.5} दिन।

\textbf{13.} संयुक्त नियम \qty{7.53e-3}{m.s^{-1}} देता है: कण \qty{8}{m} गिरता है
\qty{1062}{s} में, लगभग \qty{18}{min}। वह फँसता है।

\textbf{14.} $w_s = \num{8.99e5} \times (\num{2e-5})^2 =
\qty{3.60e-4}{m.s^{-1}}$: \qty{8}{m} \qty{22200}{s} में, लगभग \num{6.2}~घंटे। फँसता है।

\textbf{15.} $w_s = \num{8.99e5} \times (\num{1e-6})^2 =
\qty{8.99e-7}{m.s^{-1}}$: \qty{8}{m} \qty{8.9e6}{s} लेता है, \num{103}~दिन, आवास-काल का सात गुना। अधिकांश पानी के साथ बाँध से गुज़र जाता है।

\textbf{16.} कण को \qty{8}{m} \qty{1.25e6}{s} में गिरना होगा, कम से कम
$\qty{6.4e-6}{m.s^{-1}}$ पर। स्टोक्स के नियम से $d = \sqrt{18\nu
w_s/(Rg)} = \sqrt{\num{6.4e-6}/\num{8.99e5}} = \qty{2.7}{\micro m}$: एक बारीक
मृत्तिका। मोटे कण जहाँ भी प्रवेश करते हैं वहीं फँसते हैं।

\textbf{17.} $568 + 0.7 \times 5680 = \qty{4544}{t}$ प्रति वर्ष, जो
$\num{4.544e6}/1300 = \qty{3496}{m^3}$ प्रति वर्ष भरता है, जलाशय का \qty{0.233}{\percent}। उसका आधा, \qty{7.5e5}{m^3}, लगभग \num{215} वर्षों में भरता है।

\textbf{18.} $100/43.0 = \qty{2.33}{Myr}$।

\textbf{19.} $1 - 2700/3300 = 0.18$: सतह केवल हटाई गई मोटाई के
\qty{18}{\percent} तक नीचे आती है।

\textbf{20.} $100/0.1818 = \qty{550}{m}$ चट्टान।

\textbf{21.} पूरा क्षेत्र $550 \times 2700/3300 = \qty{450}{m}$ प्रत्यावर्तन से उठता है। लगभग अपरदित न हुए शिखर उसके साथ \qty{450}{m} तक ऊपर उठाए जाते हैं, जबकि घाटियाँ गहरी होती हैं: उच्चावच घटने से पहले बढ़ता है।

\textbf{22.} \omterm{def:b1:erosion-transport:denudation}{अपक्षरण-दर} सबसे अधिक ढलानों पर निर्भर करती है। जैसे-जैसे पहाड़ियाँ नीची होती हैं और उनकी ढलानें ढीली पड़ती हैं, दर घटेगी, जब तक उत्थान या जलवायु का बदलना उन्हें फिर से तीखा न करे: आज की दर से परिकलित समय एक निचली सीमा है।

\textbf{23.} $550/42.95 = 12.80$। आज की दर पर और समस्थितिक प्रत्यावर्तन के साथ, वार्डन हिल्स को \qty{100}{m} नीचे लाने में \textbf{12.8}~दस लाख वर्ष लगते हैं।
\end{solution}
